Nesta seção se aborda o tema do trabalho, a problemática da pesquisa, as
hipóteses, a justificativa, o objetivo geral e os específicos, a importância
da pesquisa para a comunidade científica e/ou sociedade, a metodologia de
pesquisa utilizada e uma rápida explanação do que trata cada seção da
monografia. A paginação da monografia inicia-se a partir desta página. A
introdução é uma seção primária.

Segue abaixo as normas da ABNT que são indispensáveis para a apresentação de
trabalhos científicos, de acordo com ABNT NBR 14724:

\begin{itemize}
  \item ABNT NBR 6023, Informação e documentação -- Referências -- Elaboração;
  \item ABNT NBR 6024, Informação e documentação -- Numeração progressiva das
        seções de um documento escrito -- Apresentação;
  \item ABNT NBR 6027, Informação e documentação -- Sumário -- Apresentação;
  \item ABNT NBR 6028, Informação e documentação -- Resumo -- Procedimento;
  \item ABNT NBR 6034, Informação e documentação -- Índice -- Apresentação;
  \item ABNT NBR 10520, Informação e documentação -- Citações em documentos --
        Apresentação;
  \item IBGE. Normas de apresentação tabular. 3. ed. Rio de Janeiro: IBGE, 1993.
        Disponível em:
        \url{https://biblioteca.ibge.gov.br/visualizacao/livros/liv23907.pdf}.
        Acesso em: 1 fev. 2022.
\end{itemize}

O texto escrito pelo autor ou as citações diretas até 3 linhas devem ser
digitado com espaçamento 1,5 entre as linhas, exceto se for citação direta com
mais de três linhas, notas de rodapé, referências, legendas das ilustrações e
das tabelas, natureza (tipo do trabalho, objetivo, nome da instituição a que é
submetido e área de concentração), que devem ser digitados em espaço simples.

\textopadrao
